\expandafter\gdef\csname SNside1\endcsname{\node[anchor=north west,inner sep=0pt,text width=48mm,font=\SideFont,align=left] at ([xshift=138mm,yshift=-466.85pt]current page.north west) {\textcolor{SNink}{\textbf{错项排除 · 有穷不等于很快}}\par\vspace{1.5mm}执行2的n次方步，只要对每个合法有限n都能结束，仍可满足有穷性。可行性也不是``考试规模下跑得快''；先判基本特性，再评价实际效率，不要混用两个层次。};}
\expandafter\gdef\csname SNside2\endcsname{\node[anchor=north west,inner sep=0pt,text width=48mm,font=\SideFont,align=left] at ([xshift=138mm,yshift=-348.56pt]current page.north west) {\textcolor{SNink}{\textbf{错项排除 · 上界可以不紧}}\par\vspace{1.5mm}线性时间当然满足O(n²)，但选择题通常要求紧数量级。若求和结果是1+2+4+\ensuremath{\cdots}，用最大项乘项数会得到合法却不紧的上界；还要检查能否消去多余的对数因子。};}
\expandafter\gdef\csname SNside3\endcsname{\node[anchor=north west,inner sep=0pt,text width=48mm,font=\SideFont,align=left] at ([xshift=138mm,yshift=-95.00pt]current page.north west) {\textcolor{SNink}{\textbf{不变量 · 累计值与轮数不同}}\par\vspace{1.5mm}第k轮结束时sum=k(k+1)/2，循环次数是满足停止不等式的最小k。sum约等于n不代表循环n次；每轮增量逐步变大，正是平方根阶出现的原因。};}
\expandafter\gdef\csname SNside4\endcsname{\node[anchor=north west,inner sep=0pt,text width=48mm,font=\SideFont,align=left] at ([xshift=138mm,yshift=-95.00pt]current page.north west) {\textcolor{SNink}{\textbf{不变量 · 先列每轮工作量}}\par\vspace{1.5mm}外层i取1、2、4等值，内层工作量恰与i相同。用等比和可知所有小轮之和仍受末轮的常数倍控制；``两层循环''没有提供平方阶的充分条件。};}
\expandafter\gdef\csname SNside5\endcsname{\node[anchor=north west,inner sep=0pt,text width=48mm,font=\SideFont,align=left] at ([xshift=138mm,yshift=-95.00pt]current page.north west) {\textcolor{SNink}{\textbf{课程连接 · 同一关系，多种表示}}\par\vspace{1.5mm}操作系统的空闲块可用位图或链表记录，表达的都是``哪些块可用''。更换表示会改变定位和更新代价。复习时把逻辑关系与地址表示分开，才能解释同一功能为何有不同算法。};}
\expandafter\gdef\csname SNside6\endcsname{\node[anchor=north west,inner sep=0pt,text width=48mm,font=\SideFont,align=left] at ([xshift=138mm,yshift=-95.00pt]current page.north west) {\textcolor{SNink}{\textbf{最小反例 · 链式不等于树}}\par\vspace{1.5mm}三个结点a\ensuremath{\to}b\ensuremath{\to}c\ensuremath{\to}a可以构成循环单链表的物理链接，链接图含环，不能据此认定为树。采用指针联系只是存储表示；逻辑对象究竟是线性表、树还是图，还需查看关系与操作约定。\par\smallskip\begin{tikzpicture}[x=1mm,y=1mm,>=Stealth,node distance=5mm]
\node[draw=SNink!50,fill=SNblue,rounded corners=1mm,text width=34mm,align=center,inner sep=1.2mm,font=\fontsize{7.5}{9}\selectfont] (m0) at (0,0) {a};
\node[draw=SNink!50,fill=SNblue,rounded corners=1mm,text width=34mm,align=center,inner sep=1.2mm,font=\fontsize{7.5}{9}\selectfont] (m1) at (0,-11) {b};
\node[draw=SNink!50,fill=SNblue,rounded corners=1mm,text width=34mm,align=center,inner sep=1.2mm,font=\fontsize{7.5}{9}\selectfont] (m2) at (0,-22) {c};
\draw[->,SNink] (m0) -- (m1);
\draw[->,SNink] (m1) -- (m2);
\draw[->,SNink] (m2) edge[bend left=40] (m0);
\end{tikzpicture}
\par{\fontsize{7.5}{9}\selectfont 循环链接含环，链式存储不能确定树关系}};}
\expandafter\gdef\csname SNside7\endcsname{\node[anchor=north west,inner sep=0pt,text width=48mm,font=\SideFont,align=left] at ([xshift=138mm,yshift=-95.00pt]current page.north west) {\textcolor{SNink}{\textbf{不变量 · 已处理区的最大值}}\par\vspace{1.5mm}进入第i轮前，best只承诺等于前i个元素的最大值，并没有承诺覆盖未扫描区。证明每轮将承诺扩展一个元素，比直接说``循环完一定正确''多出了保持关系这一关键环节。};}
\expandafter\gdef\csname SNside8\endcsname{\node[anchor=north west,inner sep=0pt,text width=48mm,font=\SideFont,align=left] at ([xshift=138mm,yshift=-95.00pt]current page.north west) {\textcolor{SNink}{\textbf{错项排除 · 效率不是唯一质量}}\par\vspace{1.5mm}一个输出错误的程序即使时间和空间都很少，也没有解决目标问题。遇到用效率``全面衡量优劣''的说法，应先检查正确性等前提；正确性、可读性、健壮性不能被遗漏。};}
\expandafter\gdef\csname SNside9\endcsname{\node[anchor=north west,inner sep=0pt,text width=48mm,font=\SideFont,align=left] at ([xshift=138mm,yshift=-95.00pt]current page.north west) {\textcolor{SNink}{\textbf{最小反例 · 零不能当通用初值}}\par\vspace{1.5mm}数组只有一个元素\ensuremath{-}3时，若best初始化为0并只在出现更大元素时修改，最后返回0就出错。非空数组应从真实元素初始化；空数组另行定义失败结果。};}
\AddToShipoutPictureFG{\begin{tikzpicture}[remember picture,overlay]\ifcsname SNside\thepage\endcsname\csname SNside\thepage\endcsname\fi\end{tikzpicture}}